% GENERATED FILE. Edit canonical lessons, then run npm run generate:books.
% canonical-source-hash: fnv1a64:b589761b2abfab22
% canonical-lessons: GU-C236-pankhi, GU-C236-limdo, GU-C236-dali, GU-C236-cac, GU-C236-canvu, GU-R236-review-231-236

\chapter{Birds in the neem tree}
\label{ch:gu-birds-in-the-neem-tree}
\hlchaptermodality{236}

\begin{chapteropening}
\textbf{By the end of this chapter:} \emph{I can say what the birds did in the tree.}

Complete the last lesson of chapter 236: I can say what the birds did in the tree.
\end{chapteropening}

\section[\texorpdfstring{\gu{પંખી} (pankhī)}{પંખી (pankhī)}]{\gu{પંખી} (pankhī) — bird}

\label{lesson:GU-C236-pankhi}



\begin{quote}
Before the new one: say the Gujarati for to lick, then the Gujarati for tail.
\end{quote}

\subsection*{You'll want to know: \gu{પંખી}}
\textbf{\gu{પંખી}} — \emph{pankhī} — \textquotedblleft{}bird\textquotedblright{}.

\begin{quote}
\textbf{\gu{સવારે} \gu{એક} \gu{પંખી} \gu{આવ્યું}.} — \emph{savāre ek pankhī āvyũ} — \textquotedblleft{}A bird came in the morning.\textquotedblright{}
\end{quote}

Not \textbf{\gu{પંખો}}, a fan.

\begin{cousinweb}
\textbf{\gu{પંખી}} (\emph{pankhī}) is a tadbhava. It comes from Sanskrit \emph{pakṣin}, \textquotedblleft{}winged; bird\textquotedblright{}, built on \emph{pakṣa}, \textquotedblleft{}wing\textquotedblright{} (CDIAL 7636).
\end{cousinweb}

\begin{grammarlens}[title={What you've built}]
One more word for everyday life at A2.
\end{grammarlens}

\subsection*{Your turn}
\begin{itemize}
\raggedright
  \item \emph{Say it:} \emph{pankhī}
  \item \emph{Say it:} \emph{pankhī}, once more
  \item \emph{Say it:} \emph{savāre ek pankhī āvyũ} — A bird came in the morning.
  \item \emph{Recall:} say the Gujarati for to lick, then the Gujarati for tail, then say \emph{pankhī} again
\end{itemize}

\begin{tcolorbox}[breakable,colback=teal!4,colframe=teal!35!black,title={Before you move on}]
What does \gu{પંખી} mean? (\textquotedblleft{}Bird\textquotedblright{}.) Say it once more.
\end{tcolorbox}
\section[\texorpdfstring{\gu{લીમડો} (līmḍo)}{લીમડો (līmḍo)}]{\gu{લીમડો} (līmḍo) — neem tree}

\label{lesson:GU-C236-limdo}



\begin{quote}
Before the new one: say the Gujarati for tail, then the Gujarati for bird.
\end{quote}

\subsection*{You'll want to know: \gu{લીમડો}}
\textbf{\gu{લીમડો}} — \emph{līmḍo} — \textquotedblleft{}neem tree\textquotedblright{}.

\begin{quote}
\textbf{\gu{અમારા} \gu{ઘર} \gu{પાસે} \gu{લીમડો} \gu{હતો}.} — \emph{amārā ghar pāse līmḍo hato} — \textquotedblleft{}There was a neem tree near our house.\textquotedblright{}
\end{quote}

\begin{cousinweb}
\textbf{\gu{લીમડો}} (\emph{līmḍo}) is a tadbhava, a word inherited from Sanskrit and reshaped by sound change. It comes from \emph{nimba}, \textquotedblleft{}the neem tree\textquotedblright{}, where CDIAL lists Gujarati \emph{lĩbṛo} (CDIAL 7245).
\end{cousinweb}

\begin{grammarlens}[title={What you've built}]
One more word for everyday life at A2.
\end{grammarlens}

\subsection*{Your turn}
\begin{itemize}
\raggedright
  \item \emph{Say it:} \emph{līmḍo}
  \item \emph{Say it:} \emph{līmḍo}, once more
  \item \emph{Say it:} \emph{amārā ghar pāse līmḍo hato} — There was a neem tree near our house.
  \item \emph{Recall:} say the Gujarati for tail, then the Gujarati for bird, then say \emph{līmḍo} again
\end{itemize}

\begin{tcolorbox}[breakable,colback=teal!4,colframe=teal!35!black,title={Before you move on}]
What does \gu{લીમડો} mean? (\textquotedblleft{}Neem tree\textquotedblright{}.) Say it once more.
\end{tcolorbox}
\section[\texorpdfstring{\gu{ડાળી} (ḍāḷī)}{ડાળી (ḍāḷī)}]{\gu{ડાળી} (ḍāḷī) — branch, twig}

\label{lesson:GU-C236-dali}



\begin{quote}
Before the new one: say the Gujarati for bird, then the Gujarati for neem tree.
\end{quote}

\subsection*{You'll want to know: \gu{ડાળી}}
\textbf{\gu{ડાળી}} — \emph{ḍāḷī} — \textquotedblleft{}branch, twig\textquotedblright{}.

\begin{quote}
\textbf{\gu{ગઈકાલે} \gu{લીમડાની} \gu{ડાળી} \gu{પડી}.} — \emph{gaīkāle līmḍānī ḍāḷī paḍī} — \textquotedblleft{}Yesterday a branch of the neem tree fell.\textquotedblright{}
\end{quote}

\begin{cousinweb}
\textbf{\gu{ડાળી}} (\emph{ḍāḷī}) is inherited from the late Sanskrit \emph{ḍāla}, \textquotedblleft{}branch\textquotedblright{}; CDIAL lists Gujarati \emph{ḍāḷī}, \textquotedblleft{}branch\textquotedblright{}, and notes that the many shapes of the word point to a non-Aryan, perhaps Munda, source (CDIAL 5546).
\end{cousinweb}

\begin{grammarlens}[title={What you've built}]
One more word for everyday life at A2.
\end{grammarlens}

\subsection*{Your turn}
\begin{itemize}
\raggedright
  \item \emph{Say it:} \emph{ḍāḷī}
  \item \emph{Say it:} \emph{ḍāḷī}, once more
  \item \emph{Say it:} \emph{gaīkāle līmḍānī ḍāḷī paḍī} — Yesterday a branch of the neem tree fell.
  \item \emph{Recall:} say the Gujarati for bird, then the Gujarati for neem tree, then say \emph{ḍāḷī} again
\end{itemize}

\begin{tcolorbox}[breakable,colback=teal!4,colframe=teal!35!black,title={Before you move on}]
What does \gu{ડાળી} mean? (\textquotedblleft{}Branch, twig\textquotedblright{}.) Say it once more.
\end{tcolorbox}
\section[\texorpdfstring{\gu{ચાંચ} (cā̃c)}{ચાંચ (cā̃c)}]{\gu{ચાંચ} (cā̃c) — beak}

\label{lesson:GU-C236-cac}



\begin{quote}
Before the new one: say the Gujarati for neem tree, then the Gujarati for branch.
\end{quote}

\subsection*{You'll want to know: \gu{ચાંચ}}
\textbf{\gu{ચાંચ}} — \emph{cā̃c} — \textquotedblleft{}beak\textquotedblright{}.

\begin{quote}
\textbf{\gu{પોપટની} \gu{ચાંચ} \gu{લાલ} \gu{હતી}.} — \emph{popaṭnī cā̃c lāl hatī} — \textquotedblleft{}The parrot's beak was red.\textquotedblright{}
\end{quote}

\begin{cousinweb}
\textbf{\gu{ચાંચ}} (\emph{cā̃c}) is inherited from Sanskrit \emph{cañcu}, \textquotedblleft{}beak\textquotedblright{}, a word CDIAL takes from a Dravidian source; it lists Gujarati \emph{cā̃c}, \textquotedblleft{}beak\textquotedblright{} (CDIAL 4569).
\end{cousinweb}

\begin{grammarlens}[title={What you've built}]
One more word for everyday life at A2.
\end{grammarlens}

\subsection*{Your turn}
\begin{itemize}
\raggedright
  \item \emph{Say it:} \emph{cā̃c}
  \item \emph{Say it:} \emph{cā̃c}, once more
  \item \emph{Say it:} \emph{popaṭnī cā̃c lāl hatī} — The parrot's beak was red.
  \item \emph{Recall:} say the Gujarati for neem tree, then the Gujarati for branch, then say \emph{cā̃c} again
\end{itemize}

\begin{tcolorbox}[breakable,colback=teal!4,colframe=teal!35!black,title={Before you move on}]
What does \gu{ચાંચ} mean? (\textquotedblleft{}Beak\textquotedblright{}.) Say it once more.
\end{tcolorbox}
\section[\texorpdfstring{\gu{ચણવું} (caṇvũ)}{ચણવું (caṇvũ)}]{\gu{ચણવું} (caṇvũ) — to peck}

\label{lesson:GU-C236-canvu}



\begin{quote}
Before the new one: say the Gujarati for branch, then the Gujarati for beak.
\end{quote}

\subsection*{You'll want to know: \gu{ચણવું}}
\textbf{\gu{ચણવું}} — \emph{caṇvũ} — \textquotedblleft{}to peck\textquotedblright{}.

\begin{quote}
\textbf{\gu{ચકલી} \gu{આંગણામાં} \gu{ચણી}.} — \emph{caklī ā̃gaṇāmā̃ caṇī} — \textquotedblleft{}The sparrow pecked in the courtyard.\textquotedblright{}
\end{quote}

\begin{cousinweb}
The origin of \textbf{\gu{ચણવું}} (\emph{caṇvũ}) in the sense \textquotedblleft{}to peck\textquotedblright{} is not settled. CDIAL lists Gujarati \emph{caṇvũ} only as \textquotedblleft{}to build\textquotedblright{}, under Sanskrit \emph{cinoti}, \textquotedblleft{}heaps up\textquotedblright{} (CDIAL 4814), and gives \emph{cugvũ} as the word for a bird pecking (CDIAL 4852).
\end{cousinweb}

\begin{grammarlens}[title={What you've built}]
One more word for everyday life at A2.
\end{grammarlens}

\subsection*{Your turn}
\begin{itemize}
\raggedright
  \item \emph{Say it:} \emph{caṇvũ}
  \item \emph{Say it:} \emph{caṇvũ}, once more
  \item \emph{Say it:} \emph{caklī ā̃gaṇāmā̃ caṇī} — The sparrow pecked in the courtyard.
  \item \emph{Recall:} say the Gujarati for branch, then the Gujarati for beak, then say \emph{caṇvũ} again
\end{itemize}

\begin{tcolorbox}[breakable,colback=teal!4,colframe=teal!35!black,title={Before you move on}]
What does \gu{ચણવું} mean? (\textquotedblleft{}To peck\textquotedblright{}.) Say it once more.
\end{tcolorbox}
\section[(dialogue)]{Words from chapters 231–236 — a review}

\label{lesson:GU-R236-review-231-236}



\begin{quote}
Say the Gujarati for pill and ointment.
\end{quote}

\subsection*{Your turn}
Say these aloud:

\begin{itemize}
\raggedright
  \item pill, ointment, nurse, treatment, to be cured
  \item disease, infection, healthy, to take care, to get well
  \item animal, horse, camel, ox, to plough
  \item calf, to moo, to milk, sheep, to graze
  \item to bark, to bite, to keep a pet, to lick, tail
  \item bird, neem tree, branch, beak, to peck
\end{itemize}

\begin{tcolorbox}[breakable,colback=teal!4,colframe=teal!35!black,title={Before you move on}]
Pill? (\textbf{\gu{ગોળી}}.) Ointment? (\textbf{\gu{મલમ}}.) Nurse? (\textbf{\gu{નર્સ}}.) Treatment? (\textbf{\gu{સારવાર}}.) To be cured? (\textbf{\gu{મટવું}}.) Disease? (\textbf{\gu{રોગ}}.) Infection? (\textbf{\gu{ચેપ}}.) Healthy? (\textbf{\gu{તંદુરસ્ત}}.) To take care? (\textbf{\gu{કાળજી} \gu{લેવી}}.) To get well? (\textbf{\gu{સાજું} \gu{થવું}}.) Animal? (\textbf{\gu{પ્રાણી}}.) Horse? (\textbf{\gu{ઘોડો}}.) Camel? (\textbf{\gu{ઊંટ}}.) Ox? (\textbf{\gu{બળદ}}.) To plough? (\textbf{\gu{ખેડવું}}.) Calf? (\textbf{\gu{વાછરડું}}.) To moo? (\textbf{\gu{ભાંભરવું}}.) To milk? (\textbf{\gu{દોહવું}}.) Sheep? (\textbf{\gu{ઘેટું}}.) To graze? (\textbf{\gu{ચરવું}}.) To bark? (\textbf{\gu{ભસવું}}.) To bite? (\textbf{\gu{કરડવું}}.) To keep a pet? (\textbf{\gu{પાળવું}}.) To lick? (\textbf{\gu{ચાટવું}}.) Tail? (\textbf{\gu{પૂંછડી}}.) Bird? (\textbf{\gu{પંખી}}.) Neem tree? (\textbf{\gu{લીમડો}}.) Branch? (\textbf{\gu{ડાળી}}.) Beak? (\textbf{\gu{ચાંચ}}.) To peck? (\textbf{\gu{ચણવું}}.)
\end{tcolorbox}
